\documentclass[11pt,a4paper]{article}
\usepackage[utf8]{inputenc}
\usepackage[vietnamese]{babel}
\usepackage{amsmath,amssymb}
\usepackage{geometry}
\geometry{a4paper, total={180mm,260mm}, left=15mm, right=15mm, top=18mm, bottom=18mm}
\usepackage{tikz}
\usetikzlibrary{shapes.geometric, arrows.meta, positioning, calc}
\usepackage{tcolorbox}
\tcbuselibrary{skins,breakable}
\usepackage{xcolor}
\usepackage{multicol}
\usepackage{array}
\usepackage{fontawesome}
\usepackage{fancyhdr}

\pagestyle{fancy}
\fancyhf{}
\renewcommand{\headrulewidth}{0.5pt}
\renewcommand{\footrulewidth}{0.5pt}
\fancyhead[L]{\small\bfseries\color{teal}\faBook\ SỔ TAY KIẾN THỨC TOÁN 3}
\fancyhead[R]{\small\bfseries\color{teal}Lớp Toán Thầy Lĩnh -- loptoanthaylinh.vn}
\fancyfoot[C]{\small\color{gray}-- Sổ tay kiến thức: Bài 42. Ôn tập biểu thức số --}

\definecolor{primary}{HTML}{005F73}
\definecolor{secondary}{HTML}{0A9396}
\definecolor{accent}{HTML}{EE9B00}
\definecolor{danger}{HTML}{AE2012}
\definecolor{boxbg}{HTML}{F4F6F8}

\begin{document}

% --- BANNER TIÊU ĐỀ ---
\begin{tcolorbox}[
  enhanced,
  colback=white,
  colframe=primary,
  arc=3mm,
  boxrule=1.2pt,
  left=4mm, right=4mm, top=2.5mm, bottom=2.5mm,
  drop shadow=black!10!white
]
  \centering
  {\Large\bfseries\color{primary}\faBookmark\ BÀI 42: ÔN TẬP BIỂU THỨC SỐ}\\[1.5mm]
  {\small\itshape\color{black!75}Ôn tập học kì 1: Hệ thống thứ tự phép tính, dấu ngoặc tròn và bài toán tính nhanh giá trị biểu thức trên 1 trang A4}
\end{tcolorbox}

\vspace{-2mm}


\noindent
\begin{minipage}[t]{0.49\textwidth}\vspace{0pt}
\begin{tcolorbox}[
  enhanced,
  colback=blue!4!white,
  colframe=blue!75!black,
  arc=2mm, boxrule=0.8pt,
  left=2mm, right=2mm, top=1.8mm, bottom=1.8mm,
  title={\footnotesize\bfseries\color{white}\faKey\ 1. Kim chỉ nam thứ tự thực hiện}
]
{\footnotesize
\begin{itemize}\setlength{\itemsep}{1.5pt}\setlength{\leftskip}{-2.5mm}
  \item \textbf{Cấp ưu tiên 1 (Dấu ngoặc tròn):}\\
  Luôn tính toán trong ngoặc $( )$ trước tiên.
  \item \textbf{Cấp ưu tiên 2 (Phép nhân và chia):}\\
  Thực hiện nhân, chia trước khi cộng, trừ.
  \item \textbf{Cấp ưu tiên 3 (Phép cộng và trừ):}\\
  Thực hiện sau khi đã hoàn tất nhân chia.
  \item \textbf{Quy tắc cùng cấp độ:} Biểu thức chỉ có $(+, -)$ hoặc chỉ có $(\times, :)$ thì làm từ \textbf{trái sang phải}.
\end{itemize}
}
\end{tcolorbox}
\end{minipage}
\hfill
\begin{minipage}[t]{0.49\textwidth}\vspace{0pt}
\begin{tcolorbox}[
  enhanced,
  colback=teal!4!white,
  colframe=teal!75!black,
  arc=2mm, boxrule=0.8pt,
  left=2mm, right=2mm, top=1.8mm, bottom=1.8mm,
  title={\footnotesize\bfseries\color{white}\faMagic\ 2. Mẹo tính nhanh giá trị biểu thức}
]
{\footnotesize
\begin{itemize}\setlength{\itemsep}{1.5pt}\setlength{\leftskip}{-2.5mm}
  \item \textbf{Nhóm thành số tròn chục, tròn trăm:}
  $$(a + b) + c = a + (b + c)$$
  \item \textbf{Đặt thừa số chung (Phân phối):}
  $$a \times b + a \times c = a \times (b + c)$$
  $$a \times b - a \times c = a \times (b - c)$$
  \item Biểu thức chứa thừa số $0$: Tích có một thừa số bằng $0$ thì cả tích bằng $0$.
\end{itemize}
}
\end{tcolorbox}
\end{minipage}

\vspace{-2.5mm}
% --- KHỐI 3: SƠ ĐỒ TƯ DUY TIKZ ---
\noindent
\begin{tcolorbox}[
  enhanced,
  colback=orange!4!white,
  colframe=orange!80!black,
  arc=2mm, boxrule=0.8pt,
  left=2mm, right=2mm, top=1.5mm, bottom=1.5mm,
  title={\footnotesize\bfseries\color{white}\faSitemap\ SƠ ĐỒ TƯ DUY: LƯỢC ĐỒ TỔNG KẾT BIỂU THỨC SỐ}
]
\centering
\begin{tikzpicture}[
  box/.style={draw, rounded corners=2pt, align=center, font=\scriptsize\bfseries, inner sep=2.5pt, minimum height=0.75cm},
  arr/.style={-{Stealth[length=2mm]}, thick}
]
  \node[box, fill=purple!15, draw=purple!80!black, text width=2.6cm] (m1) at (0, 0) {BƯỚC 1: XÉT NGOẶC\\Có ngoặc $( )$\\$\implies$ Tính trong ngoặc};
  \node[box, fill=teal!15, draw=teal!80!black, text width=2.6cm] (m2) at (3.5, 0) {BƯỚC 2: NHÂN / CHIA\\Ưu tiên cao hơn\\$\times$ và $:$ trước};
  \node[box, fill=blue!12, draw=blue!80!black, text width=2.6cm] (m3) at (7.0, 0) {BƯỚC 3: CỘNG / TRỪ\\Ưu tiên thấp hơn\\$+$ và $-$ sau};
  \node[box, fill=orange!15, draw=orange!80!black, text width=2.6cm] (m4) at (10.5, 0) {CÙNG CẤP ĐỘ\\Lần lượt từ\\TRÁI QUA PHẢI};

  \draw[arr, purple!80!black] (m1) -- (m2);
  \draw[arr, teal!80!black] (m2) -- (m3);
  \draw[arr, blue!80!black] (m3) -- (m4);
\end{tikzpicture}
\end{tcolorbox}

\vspace{-2.5mm}
% --- KHỐI 4: VÍ DỤ & MẸO TRÁNH BẪY ---
\noindent
\begin{minipage}[t]{0.49\textwidth}\vspace{0pt}
\begin{tcolorbox}[
  enhanced,
  colback=green!4!white,
  colframe=green!60!black,
  arc=2mm, boxrule=0.8pt,
  left=2mm, right=2mm, top=1.8mm, bottom=1.8mm,
  title={\footnotesize\bfseries\color{white}\faCheckCircle\ Ví dụ thực hành chuẩn mực}
]
{\footnotesize
\textbf{1. Tính thông thường:}\\
$A = 240 - 40 \times 3 = 240 - 120 = \mathbf{120}$.\\
$B = (240 - 40) \times 3 = 200 \times 3 = \mathbf{600}$.\\
\textbf{2. Tính nhanh thông minh:}\\
$C = 17 \times 4 + 17 \times 6 = 17 \times (4 + 6) = 17 \times 10 = \mathbf{170}$.\\
$D = (125 - 25 \times 5) \times 999 = (125 - 125) \times 999 = 0 \times 999 = \mathbf{0}$.
}
\end{tcolorbox}
\end{minipage}
\hfill
\begin{minipage}[t]{0.49\textwidth}\vspace{0pt}
\begin{tcolorbox}[
  enhanced,
  colback=red!4!white,
  colframe=red!70!black,
  arc=2mm, boxrule=0.8pt,
  left=2mm, right=2mm, top=1.8mm, bottom=1.8mm,
  title={\footnotesize\bfseries\color{white}\faExclamationTriangle\ Mẹo tránh bẫy trắc nghiệm}
]
{\footnotesize
\begin{itemize}\setlength{\itemsep}{1.5pt}\setlength{\leftskip}{-2.5mm}
  \item \textbf{Bẫy tính từ trái sang phải mà bỏ qua nhân chia:}\\
  Biểu thức $240 - 40 \times 3$, học sinh hay lấy $240 - 40 = 200$ rồi nhân $3$ bằng $600$! Đây là lỗi sai kinh điển, phải nhân $40 \times 3 = 120$ trước!
  \item \textbf{Bẫy bỏ quên ngoặc:}\\
  Muốn tính cộng trừ trước thì bắt buộc phải đặt dấu ngoặc $( )$.
\end{itemize}
}
\end{tcolorbox}
\end{minipage}


\end{document}
